% ======================================================================
% CAPÍTULO 4: DISEÑO E IMPLEMENTACIÓN DEL FRAMEWORK FATIGUESET-LIB (EXPANDIDO)
% ======================================================================

\chapter{Diseño e Implementación del Framework \code{fatigueset-lib}}
\label{cap:arquitectura_framework}

\section{Metodología de Ingeniería del Software y Principios de Diseño}
\label{sec:soft_principios}

El desarrollo de software en proyectos de investigación científica en Inteligencia Artificial requiere un rigor metodológico equivalente al empleado en la ingeniería de sistemas críticos \cite{sommerville2011software, guillen2025libro}. Con frecuencia, el código de experimentación en Deep Learning sufre de acoplamiento espurio, cuadernos monolíticos sin modularizar y escasa cobertura de pruebas, lo que compromete gravemente la reproducibilidad de los resultados.

Para solventar estas carencias, en este Trabajo de Fin de Grado se ha diseñado y programado desde cero el framework modular \code{fatigueset-lib} (\code{fatigueset}), apoyándose en el ecosistema científico estándar de Python con PyTorch \cite{paszke2019pytorch} y Scikit-Learn \cite{pedregosa2011scikit}, y estructurado bajo las siguientes directrices de ingeniería del software:

\begin{itemize}
    \item \textbf{Principios SOLID y Modularidad Estricta:} Cada módulo del framework encapsula una responsabilidad única e independiente. Dentro del paquete \code{fatigueset}, los submódulos de datos (\code{data}), ingeniería de características (\code{features}), redes neuronales (\code{models}), optimizadores (\code{optimizers}) y evaluación métrica (\code{evaluation}) interactúan exclusivamente mediante interfaces estables y fuertemente tipadas.
    \item \textbf{Tipado Estático Exhaustivo (\textit{Type Hinting}):} Todas las funciones, métodos y clases implementan anotaciones de tipo completas mediante el módulo \code{typing} de Python 3.10+, facilitando el análisis estático de código con herramientas como \code{mypy} y optimizando la mantenibilidad.
    \item \textbf{Desacoplamiento entre Modelos y Motores de Entrenamiento:} Las redes neuronales se definen como subclases puras de \code{torch.nn.Module}, mientras que el ciclo de vida del entrenamiento (propagación hacia adelante, cálculo de pérdidas multiobjetivo, retropropagación, recorte de gradientes, validación cruzada y \textit{early stopping}) se delega al motor centralizado \code{engine.py}.
    \item \textbf{Reproducibilidad Científica Total:} Se implementa una rutina determinista global que fija las semillas pseudoaleatorias en todos los generadores del sistema (Python \code{random}, NumPy y PyTorch CUDA/cuDNN).
\end{itemize}



\section{Arquitectura Modular del Paquete \code{fatigueset}}
\label{sec:soft_paquetes}

La Figura~\ref{fig:libreria_diagrama_expandido} muestra el diagrama de paquetes y dependencias del framework \code{fatigueset-lib}.

\begin{figure}[htbp]
    \centering
    \begin{tikzpicture}[node distance=1.1cm, auto,
        pack/.style={rectangle, draw=etsiitblue, fill=blue!8, rounded corners=6pt, text width=5.6cm, minimum height=2.3cm},
        subitem/.style={rectangle, draw=gray!40, fill=white, rounded corners=3pt, text width=5.2cm, font=\scriptsize},
        line/.style={draw=darkblue, -latex', thick}]
        
        \node [pack] (data) {\textbf{fatigueset.data}\\
        \vspace{0.1cm}
        \tikz{\node[subitem] {$\bullet$ \code{loader.py}: Ingesta CSV/XLSX multicanal\\$\bullet$ \code{dataset.py}: Dataset PyTorch y DataLoader\\$\bullet$ \code{metadata.py}: Procesamiento encuestas SSS/VAS};}
        };
        
        \node [pack, right=0.6cm of data] (feat) {\textbf{fatigueset.features}\\
        \vspace{0.1cm}
        \tikz{\node[subitem] {$\bullet$ \code{filters.py}: Butterworth IIR y Notch 50 Hz\\$\bullet$ \code{windows.py}: Sliding Windows ($W=60\,\text{s}$)\\$\bullet$ \code{extraction.py}: PSD Welch, Time, SampEn};}
        };
        
        \node [pack, below=0.9cm of data] (models) {\textbf{fatigueset.models}\\
        \vspace{0.1cm}
        \tikz{\node[subitem] {$\bullet$ \code{lstm.py}, \code{gru.py}, \code{rnn.py}\\$\bullet$ \code{tcn.py}, \code{cnn\_lstm.py}, \code{patchtst.py}\\$\bullet$ \code{xlstm.py}: Celda sLSTM nativa PyTorch\\$\bullet$ \code{foundation/}: MOMENT, Chronos, TimesFM};}
        };
        
        \node [pack, right=0.6cm of models] (opt) {\textbf{fatigueset.optimizers / eval}\\
        \vspace{0.1cm}
        \tikz{\node[subitem] {$\bullet$ \code{engine.py}: Train/Val loops, GroupKFold CV\\$\bullet$ \code{optimizers.py}: Factoría dinámica Optuna\\$\bullet$ \code{metrics.py}: $R^2$, MAE, RMSE, Pearson $r$};}
        };
        
        \path [line] (data) -- (feat);
        \path [line] (feat) -- (models);
        \path [line] (models) -- (opt);
    \end{tikzpicture}
    \caption{Diagrama estructural de submódulos y dependencias de la librería \code{fatigueset-lib}.}
    \label{fig:libreria_diagrama_expandido}
\end{figure}

\subsection{Flujo de Datos y Pipeline de Ejecución de la Librería}

La Figura~\ref{fig:flujo_libreria_datos} ilustra el ciclo de vida completo de transformación de datos dentro de \code{fatigueset-lib}. El diseño desacopla las responsabilidades en capas funcionales conectadas de forma ortogonal y limpia:

\begin{figure}[htbp]
    \centering
    \begin{tikzpicture}[node distance=0.5cm, auto,
        stepbox/.style={rectangle, draw=etsiitblue, fill=blue!6, rounded corners=4pt, text width=2.9cm, text centered, minimum height=1.5cm, font=\scriptsize},
        docbox/.style={rectangle, draw=ugrred, fill=red!6, rounded corners=4pt, text width=2.9cm, text centered, minimum height=1.5cm, font=\scriptsize},
        line/.style={draw=darkblue, -latex', thick}]
        
        % Fila 1 (Izquierda a Derecha)
        \node [docbox] (csv) {\textbf{Dataset Raw}\\23 Archivos CSV\\(Zephyr, Empatica,\\Muse, eSense)};
        \node [stepbox, right=0.4cm of csv] (loader) {\textbf{1. Ingesta / Valida}\\ \code{FatigueSetLoader}\\ \code{FatigueSetValidator}};
        \node [stepbox, right=0.4cm of loader] (processor) {\textbf{2. Procesamiento}\\ \code{FatigueSetProcessor}\\Remuestreo 64 Hz,\\Filtros Butterworth};
        \node [stepbox, right=0.4cm of processor] (extractor) {\textbf{3. Extracción}\\ \code{FeatureExtractor}\\Time, Welch PSD,\\SampEn, Poincaré};
        
        % Fila 2 (Derecha a Izquierda)
        \node [docbox, below=1.2cm of extractor] (ml_dataset) {\textbf{Dataset Tabular}\\\textbf{Machine Learning}\\108 Filas Agregadas\\$D = 92$ Features};
        \node [stepbox, below=1.2cm of processor] (tensors) {\textbf{4. Tensores 3D DL}\\Ventanas Deslizantes\\$(B, 3840, 23)$};
        \node [stepbox, below=1.2cm of loader] (modelfactory) {\textbf{5. Model Factory}\\ \code{create\_model()}\\11 Arquitecturas};
        \node [docbox, below=1.2cm of csv] (eval) {\textbf{6. Motor Engine}\\Validación GroupKFold\\Predicciones y Métricas};
        
        % Conexiones limpias ortogonales sin intersección de cajas
        \path [line] (csv) -- (loader);
        \path [line] (loader) -- (processor);
        \path [line] (processor) -- (extractor);
        \path [line] (extractor) -- (ml_dataset);
        \path [line] (processor) -- (tensors);
        \path [line] (ml_dataset) -- (tensors);
        \path [line] (tensors) -- (modelfactory);
        \path [line] (modelfactory) -- (eval);
    \end{tikzpicture}
    \caption{Diagrama de flujo integral de datos en \code{fatigueset-lib}: desde los ficheros brutos hasta la evaluación de modelos.}
    \label{fig:flujo_libreria_datos}
\end{figure}



\section{Implementación de las Arquitecturas en PyTorch}
\label{sec:soft_modelos}

A continuación, se documentan las implementaciones de las arquitecturas neuronales desarrolladas:

\subsection{Implementación de la Red Convolucional Temporal (TCN)}

El Listado~\ref{lst:tcn_code} ilustra la implementación de la arquitectura TCN en \code{fatigueset.models.tcn}. Cada bloque residual temporal (\code{TemporalBlock}) aplica convoluciones causales dilatadas con normalización por peso y capas de recorte (\code{Chomp1d}) para garantizar la causalidad temporal estricta:

\begin{lstlisting}[language=Python, caption={Implementación de la arquitectura \code{TCNRegressor} en PyTorch.}, label={lst:tcn_code}]
import torch
import torch.nn as nn
from typing import List

class Chomp1d(nn.Module):
    """Recorta los ultimos elementos temporales para preservar la causalidad."""
    def __init__(self, chomp_size: int):
        super().__init__()
        self.chomp_size = chomp_size

    def forward(self, x: torch.Tensor) -> torch.Tensor:
        return x[:, :, :-self.chomp_size].contiguous()

class TemporalBlock(nn.Module):
    """Bloque residual convolucional dilatado con Weight Normalization."""
    def __init__(self, n_inputs: int, n_outputs: int, kernel_size: int, stride: int, dilation: int, dropout: float = 0.2):
        super().__init__()
        padding = (kernel_size - 1) * dilation
        self.conv1 = nn.utils.weight_norm(nn.Conv1d(n_inputs, n_outputs, kernel_size, stride=stride, padding=padding, dilation=dilation))
        self.chomp1 = Chomp1d(padding)
        self.relu1 = nn.ReLU()
        self.dropout1 = nn.Dropout(dropout)

        self.conv2 = nn.utils.weight_norm(nn.Conv1d(n_outputs, n_outputs, kernel_size, stride=stride, padding=padding, dilation=dilation))
        self.chomp2 = Chomp1d(padding)
        self.relu2 = nn.ReLU()
        self.dropout2 = nn.Dropout(dropout)

        self.net = nn.Sequential(self.conv1, self.chomp1, self.relu1, self.dropout1,
                                 self.conv2, self.chomp2, self.relu2, self.dropout2)
        self.downsample = nn.Conv1d(n_inputs, n_outputs, 1) if n_inputs != n_outputs else None
        self.relu = nn.ReLU()

    def forward(self, x: torch.Tensor) -> torch.Tensor:
        out = self.net(x)
        res = x if self.downsample is None else self.downsample(x)
        return self.relu(out + res)

class TCNRegressor(nn.Module):
    """Regresor TCN con multiples niveles de dilatacion exponencial."""
    def __init__(self, num_inputs: int = 23, num_channels: List[int] = [32, 64, 64, 128], kernel_size: int = 3, dropout: float = 0.2, output_dim: int = 2):
        super().__init__()
        layers = []
        num_levels = len(num_channels)
        for i in range(num_levels):
            dilation_size = 2 ** i
            in_channels = num_inputs if i == 0 else num_channels[i - 1]
            out_channels = num_channels[i]
            layers.append(TemporalBlock(in_channels, out_channels, kernel_size, stride=1, dilation=dilation_size, dropout=dropout))
        self.network = nn.Sequential(*layers)
        self.head = nn.Linear(num_channels[-1], output_dim)

    def forward(self, x: torch.Tensor) -> torch.Tensor:
        x = x.permute(0, 2, 1)
        y = self.network(x)
        last_step = y[:, :, -1]
        return self.head(last_step)
\end{lstlisting}

\subsection{Implementación de la Arquitectura PatchTST}

El Listado~\ref{lst:patchtst_code} muestra la implementación de \code{PatchTSTRegressor}, que agrupa puntos temporales en parches continuos y aplica auto-atención multi-cabeza:

\begin{lstlisting}[language=Python, caption={Implementación de \code{PatchTSTRegressor} con parcheo e independencia de canales.}, label={lst:patchtst_code}]
import torch
import torch.nn as nn

class PatchTSTRegressor(nn.Module):
    """Transformer basado en parches para series temporales multicanal."""
    def __init__(
        self,
        num_channels: int = 23,
        patch_len: int = 32,
        stride: int = 8,
        seq_len: int = 3840,
        d_model: int = 64,
        nhead: int = 4,
        num_layers: int = 3,
        dim_feedforward: int = 128,
        dropout: float = 0.1,
        output_dim: int = 2
    ):
        super().__init__()
        self.num_channels = num_channels
        self.patch_len = patch_len
        self.stride = stride
        self.num_patches = (seq_len - patch_len) // stride + 1
        
        self.patch_embed = nn.Linear(patch_len, d_model)
        self.pos_encoder = nn.Parameter(torch.zeros(1, self.num_patches, d_model))
        
        encoder_layer = nn.TransformerEncoderLayer(
            d_model=d_model, nhead=nhead, dim_feedforward=dim_feedforward,
            dropout=dropout, batch_first=True, activation='gelu'
        )
        self.transformer_encoder = nn.TransformerEncoder(encoder_layer, num_layers=num_layers)
        self.head = nn.Linear(num_channels * d_model, output_dim)

    def forward(self, x: torch.Tensor) -> torch.Tensor:
        B, T, C = x.shape
        patches = x.unfold(dimension=1, size=self.patch_len, step=self.stride)
        patches = patches.permute(0, 2, 1, 3).reshape(B * C, self.num_patches, self.patch_len)
        
        embed = self.patch_embed(patches) + self.pos_encoder
        out = self.transformer_encoder(embed)
        out_pooled = out.mean(dim=1)
        out_flattened = out_pooled.view(B, C * embed.shape[-1])
        return self.head(out_flattened)
\end{lstlisting}

\subsection{Estructura Teórica y Flujo Interno de la Celda sLSTM (\xlstm{})}

Para comprender a fondo la ventaja operativa de la arquitectura \xlstm{}, la Figura~\ref{fig:esquema_slstm_cell} desglosa el flujo matemático interno de la celda \slstm{}. A diferencia de la LSTM clásica, donde las compuertas sigmoides se saturan entre $0$ y $1$, la compuerta exponencial $\exp(\tilde{i}_t)$ permite la sobrescritura dinámica completa del estado celular. El factor estabilizador $m_t = \max(\tilde{f}_t + m_{t-1}, \tilde{i}_t)$ actúa como un exponente de escala que previene desbordamientos en punto flotante (\textit{IEEE-754}), mientras que el estado normalizador $n_t$ preserva la escala unitaria del gradiente.

\begin{figure}[htbp]
    \centering
    \begin{tikzpicture}[node distance=1.0cm, auto,
        block/.style={rectangle, draw=darkblue, fill=blue!8, rounded corners=5pt, text width=3.7cm, text centered, minimum height=1.5cm, font=\small},
        line/.style={draw=darkblue, -latex', thick}]
        
        % Fila 1 (Izquierda a Derecha)
        \node [block] (input) {\textbf{Entrada y Estado}\\$\bm{x}_t \in \mathbb{R}^C,\; \bm{h}_{t-1} \in \mathbb{R}^d$};
        \node [block, right=0.8cm of input] (linear) {\textbf{Proyecciones Lineales}\\$\tilde{f}_t, \tilde{i}_t, \tilde{z}_t, \tilde{o}_t \in \mathbb{R}^d$};
        \node [block, right=0.8cm of linear] (stabilizer) {\textbf{Estabilizador}\\$m_t = \max(\tilde{f}_t + m_{t-1}, \tilde{i}_t)$};
        
        % Fila 2 (Derecha a Izquierda)
        \node [block, below=1.2cm of stabilizer] (expgates) {\textbf{Compuertas Exp.}\\$f_t = \exp(\tilde{f}_t - m_t + m_{t-1})$\\$i_t = \exp(\tilde{i}_t - m_t)$};
        \node [block, below=1.2cm of linear] (states) {\textbf{Memoria y Normalizador}\\$n_t = f_t n_{t-1} + i_t$\\$c_t = f_t c_{t-1} + i_t \tanh(\tilde{z}_t)$};
        \node [block, below=1.2cm of input] (out) {\textbf{Estado Normalizado}\\$\bm{h}_t = \sigma(\tilde{o}_t) \odot \left(\frac{c_t}{n_t + \epsilon}\right)$};
        
        \path [line] (input) -- (linear);
        \path [line] (linear) -- (stabilizer);
        \path [line] (stabilizer) -- (expgates);
        \path [line] (expgates) -- (states);
        \path [line] (states) -- (out);
    \end{tikzpicture}
    \caption{Diagrama de flujo del flujo interno de datos en la celda \slstm{} de la arquitectura \xlstm{}, destacando la normalización logarítmica y la estabilización numérica con $m_t$ y $n_t$.}
    \label{fig:esquema_slstm_cell}
\end{figure}



\section{Motor Centralizado de Entrenamiento y Validación Cruzada}
\label{sec:soft_engine}

El módulo \code{fatigueset.models.engine} centraliza la ejecución de los experimentos, garantizando la homogeneidad metodológica en todas las pruebas.

\subsection{Algoritmo de Validación Cruzada GroupKFold}
El Listado~\ref{lst:engine_code} documenta la función \code{train\_kfold\_cv}, encargada de coordinar el entrenamiento distribuido por particiones, registrar métricas por fold y exportar reportes estructurados en formato JSON:

\begin{lstlisting}[language=Python, caption={Bucle principal de validación cruzada GroupKFold en \code{engine.py}.}, label={lst:engine_code}]
def train_kfold_cv(
    model_fn,
    dataset,
    groups: np.ndarray,
    n_splits: int = 5,
    epochs: int = 50,
    lr: float = 1e-3,
    weight_decay: float = 1e-4,
    device: str = "cuda",
    early_stopping_patience: int = 10
) -> Dict[str, Any]:
    gkf = GroupKFold(n_splits=n_splits)
    fold_metrics = []
    
    for fold, (train_idx, val_idx) in enumerate(gkf.split(dataset, groups=groups)):
        train_sub = torch.utils.data.Subset(dataset, train_idx)
        val_sub = torch.utils.data.Subset(dataset, val_idx)
        
        train_loader = DataLoader(train_sub, batch_size=32, shuffle=True)
        val_loader = DataLoader(val_sub, batch_size=32, shuffle=False)
        
        model = model_fn().to(device)
        optimizer = torch.optim.AdamW(model.parameters(), lr=lr, weight_decay=weight_decay)
        loss_fn = nn.SmoothL1Loss()
        
        best_val_loss = float('inf')
        patience_counter = 0
        
        for epoch in range(epochs):
            train_loss = train_step(model, train_loader, loss_fn, optimizer, device)
            val_loss, val_mae, val_rmse, val_r2 = evaluate_step(model, val_loader, loss_fn, device)
            
            if val_loss < best_val_loss:
                best_val_loss = val_loss
                best_metrics = {'mae': val_mae, 'rmse': val_rmse, 'r2': val_r2}
                patience_counter = 0
            else:
                patience_counter += 1
                if patience_counter >= early_stopping_patience:
                    break
                    
        fold_metrics.append(best_metrics)
        
    return aggregate_cv_metrics(fold_metrics)
\end{lstlisting}



\section{Optimización Bayesiana y Almacenamiento SQLite con Optuna}
\label{sec:soft_optuna_arch}

Para garantizar la reproducibilidad y la trazabilidad de los experimentos a gran escala, el framework integra el motor de optimización bayesiana Optuna \cite{akiba2019optuna}, basado en el algoritmo de Estimadores Parzen Estructurados en Árbol (TPE, \textit{Tree-structured Parzen Estimator}) \cite{bergstra2011algorithms}, y conectado a una base de datos relacional SQLite persistente (\code{optuna\_v2.db}). La Figura~\ref{fig:optuna_arquitectura_slurm} esquematiza la interacción distribuida entre los workers en el clúster GPU y la base de datos centralizada.

\begin{figure}[htbp]
    \centering
    \begin{tikzpicture}[node distance=1.2cm, auto,
        worker/.style={rectangle, draw=etsiitblue, fill=blue!6, rounded corners=5pt, text width=3.6cm, text centered, minimum height=1.4cm, font=\small},
        db/.style={cylinder, draw=ugrred, fill=red!10, shape border rotate=90, aspect=0.25, text width=2.8cm, text centered, minimum height=2.0cm, font=\small\bfseries},
        optuna/.style={rectangle, draw=darkblue, fill=blue!12, rounded corners=5pt, text width=4.0cm, text centered, minimum height=1.5cm, font=\small},
        line/.style={draw=darkblue, -latex', thick}]
        
        \node [optuna] (tpe) {\textbf{Optuna TPE Sampler}\\Sugerencia Bayesiana\\de Hiperparámetros};
        \node [db, right=2.6cm of tpe] (sqlite) {Base de Datos SQLite\\\code{optuna\_v2.db}\\(Transacciones ACID)};
        
        \node [worker, below=1.4cm of tpe, xshift=-1.8cm] (w1) {\textbf{Worker SLURM 1}\\NVIDIA RTX 2080 Ti\\(Fold 1--2)};
        \node [worker, below=1.4cm of tpe, xshift=2.2cm] (w2) {\textbf{Worker SLURM 2}\\NVIDIA RTX 2080 Ti\\(Fold 3--5)};
        
        \path [line] (tpe) -- node [above=2pt, font=\scriptsize\bfseries] {Registro RDB} (sqlite);
        \path [line] (tpe.south) -| (w1.north);
        \path [line] (tpe.south) -| (w2.north);
        \path [line] (w1.south) -- ++(0,-0.4) -| node [near end, above=2pt, font=\scriptsize] {Métricas Fold} (sqlite.south);
        \path [line] (w2.south) -- ++(0,-0.4) -| (sqlite.south);
    \end{tikzpicture}
    \caption{Arquitectura distribuida de optimización bayesiana con Optuna y SQLite: coordinación de múltiples procesos SLURM en GPU con persistencia transaccional.}
    \label{fig:optuna_arquitectura_slurm}
\end{figure}

El estudio conjunto de hiperparámetros optimiza simultáneamente:
\begin{enumerate}
    \item \textbf{Hiperparámetros Estructurales:} Número de capas ($1 - 4$), dimensión oculta (\code{hidden\_size} $\in [32, 256]$) y tasa de \code{dropout} ($0.0 - 0.5$).
    \item \textbf{Optimizador y Política de Regularización:} Selección categórica entre Adam, AdamW (con decaimiento de pesos desacoplado \cite{loshchilov2018decoupled}) y RMSprop, optimizando conjuntamente la tasa de aprendizaje (\code{lr} $\in [10^{-4}, 10^{-2}]$) y el decaimiento de pesos (\code{weight\_decay} $\in [10^{-6}, 10^{-2}]$).
    \item \textbf{Poda Dinámica (\textit{Pruning}):} Implementación de \code{optuna.pruners.MedianPruner} para cancelar tempranamente configuraciones de hiperparámetros que divergen o arrojan pérdidas sustancialmente peores que la mediana histórica.
\end{enumerate}

La arquitectura modular de \code{fatigueset-lib} permite registrar y monitorizar estas trayectorias de optimización de forma sistemática. La evaluación empírica de convergencia y la comparativa de dispersión de pérdidas entre optimizadores (Adam, AdamW y RMSprop) sobre los modelos evaluados se analizan en detalle en el Capítulo~\ref{cap:resultados} (Sección~\ref{sec:res_globales}, Figura~\ref{fig:optuna_curves}).
